% !PRINTED 188
\weberpage{188}

Solche Gebilde sind die regulären Pyramiden, die Doppelpyramiden
und die regulären Körper (Tetra\"eder, Octa\"eder, Würfel,
Dodeka\"eder und Ikosa\"eder).

Die reguläre Pyramide von $g$ Seiten gelangt durch Drehung um die
Hauptaxe mit einem Winkel $2\pi:g$ und durch Wiederholung dieser
Drehung auf $g$ verschiedene Arten mit sich zur Deckung.

Ist $g$ gerade, so kann man eine reguläre Doppelpyramide von $g$
Seitenflächen ausser durch Drehung um ihre Hauptaxe mit dem Winkel
$4\pi:g$ noch (auf $\tfrac{1}{2}g$ verschiedene Arten) durch Drehung
um eine in der Aequatorialebene liegende Axe mit einem
Drehungswinkel von 180 Grad mit sich zur Deckung bringen.

Das Tetra\"eder gestattet 12 verschiedene Stellungen, in denen es
denselben Raum einnimmt.

Um sie zu erhalten, bezeichne man den Ort einer Ecke in der
Einheitsstellung mit 1. Dann erhält man drei Stellungen des
Tetra\"eders, bei denen die Ecke 1 fest bleibt. Man kann aber jede
Ecke an die Stelle von 1 bringen, wodurch die Zahl sich
vervierfacht.

Dieselbe Zahl findet man auch, wenn man die Ebene einer
Seitenfläche festhält, wobei man noch drei Stellungen des
Tetra\"eders findet, und jede Seitenfläche in die Ausgangsebene
bringt.

Endlich kann man auch so zählen, dass man jede Kante auf zwei
Arten mit einer festen Strecke zur Deckung bringt. Ebenso
verfährt man bei den übrigen regulären Körpern. Man findet so den
Grad der Gruppe

\begin{quote}
gleich dem Producte aus der Anzahl der Ecken mit der Anzahl der in
einer Ecke zusammenstossenden Kanten oder Seitenflächen, oder

gleich dem Producte aus der Anzahl der Seitenflächen mit der
Anzahl der Seiten oder Ecken einer Grenzfläche, oder

gleich der doppelten Anzahl der Kanten.
\end{quote}

Jede dieser Zählungen ergiebt

\begin{quote}
für das Tetra\"eder 12 Stellungen,\\
für das Octa\"eder und den Würfel 24 Stellungen,\\
für das Dodeka\"eder und Ikosa\"eder 60 Stellungen.
\end{quote}
